% Final report for COMP8240 Assessment Task #3.
% Uses the ACL style files, as required by the brief.
% Files required alongside this one: acl.sty, acl_natbib.bst, references.bib
% Build: pdflatex report; bibtex report; pdflatex report; pdflatex report
% Remove every \todo before submitting.

\documentclass[11pt]{article}
\usepackage[]{acl}
\usepackage{times}
\usepackage{latexsym}
\usepackage{microtype}
\usepackage{graphicx}
\usepackage{booktabs}
\usepackage{amsmath}
\usepackage{amssymb}
\usepackage{xcolor}

\newcommand{\todo}[1]{\textcolor{red}{\textbf{[TODO: #1]}}}

\title{Replicating and Extending \\ \textit{How Powerful are Graph Neural Networks?}}

\author{
  Sobhan Salarian \\
  Macquarie University \\
  COMP8240: Applications of Data Science \\
  \texttt{shaygan.sobhan79@gmail.com}
}

\begin{document}
\raggedbottom
\maketitle

\begin{abstract}
We replicate the graph classification results of \citet{xu2019gin} on IMDB-BINARY and IMDB-MULTI using the authors' code, and test the Graph Isomorphism Network (GIN) on new data. Replicated accuracy falls within one reported standard deviation of the paper on both datasets. Our first new dataset, built from The Movie Database (TMDb), produced a large drop in accuracy, which we traced to a construction flaw: every graph was a clique, leaving node count as the only signal. We rebuilt the data as actor ego-networks and also ran the existing ENZYMES benchmark. On ego-networks GIN and GraphSAGE tied at well above the node-count baseline (74.4\% and 59.7\% for GIN), while on ENZYMES GIN clearly won, 48.5\% to 33.3\%. An LLM judge agreed strongly with the genre labels ($\kappa \geq 0.946$) and its disagreements exposed a date filter bug in our scraper.
\end{abstract}

\section{Description of Paper and Justification}
\label{sec:paper}

\citet{xu2019gin}\footnote{\url{https://arxiv.org/abs/1810.00826}} ask how expressive message-passing Graph Neural Networks (GNNs) are, meaning which graphs they can tell apart. They prove that the Weisfeiler-Lehman (WL) test \citep{weisfeiler1968reduction} is an upper bound for this class of model, and that mean and max aggregation, as used in GCN \citep{kipf2017gcn} and GraphSAGE \citep{hamilton2017graphsage}, are strictly weaker. They then propose GIN, which sums neighbour features and passes them through a multi-layer perceptron (MLP), and show that it reaches the WL bound.

\textbf{Evaluation.} The paper reports graph classification accuracy on nine benchmarks using 10-fold cross-validation. For each epoch, test accuracy is averaged across the folds; the single epoch with the highest average is selected, and the mean and standard deviation across folds at that epoch are reported. The training setup is 5 GNN layers, 2-layer MLPs, 64 hidden units, batch size 32, dropout 0.5, and Adam at a learning rate of 0.01, halved every 50 epochs, for 350 epochs. GIN-0, the variant with $\epsilon$ fixed at 0, reaches 75.1 $\pm$ 5.1\% on IMDB-BINARY and 52.3 $\pm$ 2.8\% on IMDB-MULTI. The GraphSAGE approximation, Max--1-layer, reaches 72.3 $\pm$ 5.3\% and 50.9 $\pm$ 2.2\%.

\textbf{Justification.} ICLR is ranked A* by CORE. The paper has \todo{current Google Scholar citation count and access date}. Its official code is public, and its theory gives an explicit reason to expect a result to hold or fail, which suits an investigation on new data. The task also transfers well: any dataset in the same text format can be run with no change to the model.

\section{Description of the Original Dataset}
\label{sec:data}

We replicate on two of the paper's nine benchmarks, because they are the two our new data mirrors. Both were released by \citet{yanardag2015deep} and ship with the official repository.\footnote{\url{https://github.com/weihua916/powerful-gnns}} In each graph, nodes are actors and an edge means two actors appeared in a film together. Each graph is one actor's ego-network, labelled by genre: Action or Romance for IMDB-BINARY, and Comedy, Romance or Science Fiction for IMDB-MULTI. There are no node features, so node degree is used as the input feature (\texttt{--degree\_as\_tag}), which forces the model to rely on structure. Table~\ref{tab:stats} gives the statistics we computed from the files.

Each dataset is a single text file, shown in Figure~\ref{fig:snippet}. The first line is the graph count. Each graph then starts with its node count and label, followed by one line per node giving its tag, its neighbour count and its neighbours.

\begin{figure}[t]
\small
\begin{verbatim}
1000            <- number of graphs
20 0            <- 20 nodes, label 0
0 5 2 4 5 9 10  <- tag, degree, neighbours
\end{verbatim}
\caption{Start of the IMDB-BINARY text file, with our annotations added.}
\label{fig:snippet}
\end{figure}

\begin{table*}[t]
\centering
\small
\begin{tabular}{lrrrrrrr}
\toprule
Dataset & Graphs & Classes & Avg.\ nodes & Density & Cliques & Degree seqs & Node-count acc. \\
\midrule
IMDB-BINARY & 1000 & 2 & 19.8 & 0.521 & 13.9\% & 537 & 58.1\% \\
First scrape, film graphs, binary & 996 & 2 & 19.2 & 1.000 & 100\% & 17 & 56.0\% \\
Final ego-networks, binary & 1000 & 2 & 45.1 & 0.508 & 0.0\% & 336 & 57.7\% \\
\midrule
IMDB-MULTI & 1500 & 3 & 13.0 & 0.772 & 52.6\% & 387 & 40.1\% \\
First scrape, film graphs, multi & 1492 & 3 & 18.9 & 1.000 & 100\% & 18 & 40.5\% \\
Final ego-networks, multi & 1500 & 3 & 43.8 & 0.515 & 0.1\% & 633 & 46.3\% \\
\midrule
ENZYMES & 600 & 6 & 32.6 & 0.160 & 0.5\% & 586 & 23.5\% \\
\bottomrule
\end{tabular}
\caption{Dataset statistics computed from the files. Degree seqs is the number of distinct degree sequences. Node-count acc.\ is the 10-fold accuracy of a classifier that sees only the number of nodes.}
\label{tab:stats}
\end{table*}

\section{Replication of Original Work}
\label{sec:replication}

\textbf{Environment.} The code targets PyTorch 0.4.1 and 1.0, but it ran on Python 3.11 with PyTorch 2.x with only deprecation warnings. A laptop CPU needed about 3.5 hours per dataset, so we first used a Google Colab T4 GPU. When Colab's quota ran out, we moved to a university server with an RTX 6000 Ada GPU. All scripts are in our public repository.\footnote{\url{https://github.com/SobhanSalarian/Gin-Graph-application-of-data-science-}}

\textbf{Process.} The official \texttt{main.py} trains one fold, and its output file is overwritten each epoch, so it cannot hold a learning curve. We wrote a driver, \texttt{run\_gin.py}, that runs all 10 folds, parses each epoch's accuracy, and applies the paper's model-selection rule. Hyperparameters are the code's defaults, which match the paper. As in the paper, GraphSAGE is approximated in the same code with \texttt{--neighbor\_pooling\_type max --num\_mlp\_layers 1}. It is not the implementation of \citet{hamilton2017graphsage}.

\textbf{Changes to the original code.} We made one change. In the max-pooling path, \texttt{\_\_preprocess\_neighbors\_maxpool} left an index tensor on the CPU, and on GPU the first GraphSAGE fold produced no output for over 12 minutes. Adding \texttt{.to(self.device)} resolved this and does not change the computation. Separately, \texttt{scikit-learn} is missing from the repository's requirements file, so our setup script installs it.

\textbf{Results.} Table~\ref{tab:replication} compares our results with the paper. All are within one of the paper's standard deviations. The ranking also holds on IMDB-BINARY, where GIN beats GraphSAGE by 1.3 points against the paper's 2.8. On IMDB-MULTI the two differ by 0.3 points in GraphSAGE's favour, where the paper has GIN ahead by 1.4. Both gaps are well inside one standard deviation, so we treat this as a tie and do not read it as a contradiction.

A separate GIN run on Colab gave 74.4 $\pm$ 4.1\% and 50.5 $\pm$ 2.3\%, within about 0.5 points of the server. Only summary figures were kept from that run, so the table uses the server results, whose per-epoch curves are in the repository. Small differences from the paper are expected, because seeds and hardware differ. The selected epoch also varied between runs on IMDB-BINARY (13 on the server, 80 on Colab), since its averaged curve is very flat.

\begin{table}[t]
\centering
\small
\begin{tabular}{lcc}
\toprule
 & IMDB-BINARY & IMDB-MULTI \\
\midrule
Paper GIN-0 & 75.1 $\pm$ 5.1 & 52.3 $\pm$ 2.8 \\
\textbf{Ours GIN-0} & \textbf{73.9 $\pm$ 2.7} & \textbf{50.6 $\pm$ 2.7} \\
\midrule
Paper Max--1-layer & 72.3 $\pm$ 5.3 & 50.9 $\pm$ 2.2 \\
\textbf{Ours Max--1-layer} & \textbf{72.6 $\pm$ 3.0} & \textbf{50.9 $\pm$ 2.8} \\
\bottomrule
\end{tabular}
\caption{Replication on the original data: test accuracy (\%), mean $\pm$ standard deviation across 10 folds.}
\label{tab:replication}
\end{table}

\section{Construction of New Data}
\label{sec:newdata}

\subsection{Existing Dataset: ENZYMES}

ENZYMES \citep{borgwardt2005protein} contains 600 protein tertiary structures, 100 for each of 6 enzyme classes, and is not among the paper's nine benchmarks. Nodes are secondary structure elements with one of 3 labels, which we use as the input features. We downloaded it from the TUDataset collection \citep{morris2020tudataset}.\footnote{\url{https://chrsmrrs.github.io/datasets/}} Our proposal said this collection already matches the code's input format. That was wrong: it uses separate files for edges, graph membership and labels, so we wrote \texttt{convert\_tu\_dortmund.py} to produce the format in Figure~\ref{fig:snippet}.

\subsection{Constructed Dataset: TMDb}

\textbf{Collection.} We used the TMDb API \citep{tmdb} to collect films released from 2016 onward, so none overlap with the 2015 originals. For each target genre we took up to 500 films by popularity with at least 20 votes, then fetched their credits. Films tagged with more than one target genre were dropped to keep labels unambiguous. We kept the top 20 billed actors and dropped films with fewer than 5. No human annotation was needed, because labels are TMDb's genre metadata. Section~\ref{sec:judge} checks their quality. Our first scrape did not apply the date limit: it sent \texttt{primary\_release\_date\_gte} where the API expects \texttt{primary\_release\_date.gte}, and the API ignored the unknown parameter without an error, so films such as \textit{Casablanca} were included. We fixed the parameter, re-scraped, and made the scraper stop if any film falls before the cut-off. The final data has 993 binary and 1,484 multi films, dated 14 January 2016 to 29 October 2026.

\textbf{First construction, and its flaw.} In the first scrape we built one graph per film, using its cast as nodes and linking actors who co-appeared in any sampled film of that genre. Before training we did not check the resulting structure closely, and the proposal's expectation of ego-network structure turned out to be wrong. Every pair of cast members has already co-appeared in that film, so the rule adds nothing, and all 2,488 graphs were cliques (Table~\ref{tab:stats}). Two cliques of the same size are isomorphic, so node count was the only usable information.

\textbf{Rebuilt construction.} For the final data we followed \citet{yanardag2015deep} more closely. For each genre we form one collaboration graph, then take each actor's ego-network within it: the actor, their co-stars, and all links among those co-stars. Actors with only one film in a genre are excluded, because their ego-network is that film's clique. We sampled 500 graphs per genre at random from the 627 to 1,323 available.

The rebuilt data is close to the originals in density and has real structural variety (Table~\ref{tab:stats}). There are two main differences. Graphs are more than twice the size of IMDB-BINARY's, because each film contributes up to 20 actors. Cliques are almost absent because of the exclusion above, against 13.9\% and 52.6\% in the originals. As in the originals, ego-networks overlap, so train and test graphs can share actors.

\section{Results on New Data}
\label{sec:results}

Hyperparameters and protocol are unchanged, and the code is the same as in the replication. Table~\ref{tab:newresults} gives the results. The node-count column is the reference point: accuracy above it requires information from graph structure.

\begin{table}[t]
\centering
\small
\begin{tabular}{lccc}
\toprule
Dataset & GIN-0 & Max--1-layer & Node count \\
\midrule
First scrape, binary & 55.9 $\pm$ 2.3 & 56.2 $\pm$ 2.3 & 56.0 \\
First scrape, multi & 41.0 $\pm$ 2.3 & 41.0 $\pm$ 2.2 & 40.5 \\
Ego, binary & 74.4 $\pm$ 2.5 & 74.0 $\pm$ 1.9 & 57.7 \\
Ego, multi & 59.7 $\pm$ 4.4 & 57.6 $\pm$ 5.0 & 46.3 \\
ENZYMES & 48.5 $\pm$ 7.0 & 33.3 $\pm$ 5.6 & 23.5 \\
\bottomrule
\end{tabular}
\caption{Results on new data: test accuracy (\%). Node count is 10-fold accuracy using only the number of nodes.}
\label{tab:newresults}
\end{table}

\textbf{First-scrape film graphs.} GIN fell by 18.0 points on the binary task and 9.6 on the multi task compared with our replication, and tied with GraphSAGE. Both models sit at the node-count baseline, so they learned cast size and nothing else. With no structure to separate graphs, this is what the theory predicts, and it tests nothing about GIN. It does show the value of running a structure-free baseline before training.

\textbf{Ego-networks.} Unlike the first scrape, both models now sit well above the node-count baseline: GIN by 16.7 points on the binary task and 13.4 on the multi task, so they are using structure. Binary accuracy is close to our replication (74.4\% against 73.9\%). Multi accuracy is 9.1 points higher than the replication, but this does not show GIN doing better, because these graphs are larger and the baseline itself is 6.2 points higher. GIN leads GraphSAGE by 0.4 and 2.1 points, both within one standard deviation, so we treat the models as tied on this data.

\textbf{ENZYMES.} This is the only dataset with a clear gap. GIN reached 48.5\% and GraphSAGE 33.3\%, a difference of 15.2 points, more than twice either standard deviation. Both beat the node-count baseline (23.5\%) and chance (16.7\%). Mean training accuracy at the selected epoch was 99.7\% for GIN and 39.2\% for GraphSAGE, so GraphSAGE underfits while GIN overfits heavily. This fits the paper's theory: with only 3 node labels, neighbours often share a label, and max aggregation cannot count them. We did not test this cause directly, and the 1-layer MLP may also contribute.

\textbf{Limits.} The selection rule picks the epoch using test folds, as in the paper, so all figures are somewhat optimistic. Train and test ego-networks can also share actors.

\textbf{Hardware.} All runs used the GPU except GraphSAGE on the ego-networks. Max pooling was slow on these dense graphs: about 26 seconds per epoch on the shared GPU, against 13 on one CPU core. We ran those two jobs on CPU, 10 folds at a time with one thread each, which took under 2 hours per dataset rather than an estimated 25. Seeds and settings were unchanged, but their results are not numerically identical to what a GPU run would give.

\textbf{Failure analysis.} \todo{Misclassified graphs by size, density and class, if completed.}

\subsection{LLM as Judge}
\label{sec:judge}

GIN's output is a class label scored against ground truth, so accuracy needs no judge. The part of the pipeline nobody has reviewed is the labels themselves, which come from TMDb and feed every graph in our data. We therefore used an LLM as an independent annotator. This differs slightly from the proposal, which suggested checking misclassified graphs: a random sample of films tests the labels themselves, independently of the model.

We sampled 100 films per genre at random from each dataset and showed Claude Sonnet 5.5 each film's title and top 8 billed actors in shuffled batches. The TMDb label was hidden and the judge had no tools. It chose one genre or answered UNSURE, and we measured agreement with TMDb using Cohen's kappa \citep{cohen1960coefficient}, excluding UNSURE answers (Table~\ref{tab:judge}).

\begin{table}[t]
\centering
\small
\begin{tabular}{lcc}
\toprule
 & Binary & Multi \\
\midrule
Films judged & 200 & 300 \\
UNSURE & 12 & 50 \\
Agreement, all films & 94.0\% & 80.3\% \\
Agreement, decided films & 100\% & 96.4\% \\
Cohen's kappa, decided films & 1.000 & 0.946 \\
\bottomrule
\end{tabular}
\caption{LLM judge agreement with TMDb genre labels.}
\label{tab:judge}
\end{table}

Agreement is very high, which supports using TMDb's labels. The judge decided 188 binary films and agreed on all of them, which is unsurprising given how distinct Action and Romance are. Most of the 9 disagreements on the multi task involve films that plausibly fit two genres, such as \textit{Thor: Love and Thunder} (TMDb Comedy, judge Science Fiction) and \textit{Good Luck to You, Leo Grande} (TMDb Romance, judge Comedy). The multi task is harder, with 50 UNSURE answers, 24 of them for Romance films, so that label appears the least obvious from title and cast. Agreement across all films is therefore lower there (80.3\%). These figures need care. The judge chooses only among our target genres, so they do not show that a genre is a film's main one. It may also remember public genre listings, in which case it is not fully independent of TMDb, and the high figures cannot rule out shared errors. Kappa also excludes the films the judge found hardest.

The judge's most useful result was unplanned. When we first ran it on the earlier scrape, reading its disagreements exposed pre-2016 films in data meant to be 2016 onward, which led us to the date filter bug in Section~\ref{sec:newdata}. The figures above are from a rerun on the corrected films.

\section{Ethical Considerations}
\label{sec:ethics}

\todo{Rework using the unit's ethics resources, which were not available when drafting.}

\textbf{Original work.} The IMDB benchmarks were derived from public film credits, not from data about private behaviour, and the paper shares its code and data, which supports verification. However, the benchmarks describe real people. \todo{Check how the original data was collected and under what terms.}

\textbf{New data and terms of use.} The Australian Code \citep{nhmrc2018code} expects research to be honest, transparent and accountable, and to comply with relevant conditions. This matters most for our data source. TMDb's terms \citep{tmdbterms} do not permit commercial use, and they treat training or validating machine learning systems on its content as commercial use needing a written agreement. They do not mention academic use, cache data for at most six months, and require attribution. Our proposal assumed academic use was unproblematic, and we did not read these clauses before collecting. \textbf{Licence and use declaration.} The new datasets and metadata in this project were created for non-commercial university coursework only. They must not be used for any commercial purpose. Any reuse, redistribution or publication beyond this assessment requires the appropriate licence and written permission from TMDb. No such licence is claimed here, and our repository does not grant one. Publishing a research paper on this data would need the same permission.

\textbf{Attribution.} This product uses the TMDb API but is not endorsed or certified by TMDb.

\textbf{People in the data.} Our metadata holds real actors' names. These are public credits, but a genre-labelled collaboration graph can still be sensitive, because it reveals who works with whom and can reinforce associations with genres. Genre labels may also reflect the biases of TMDb's contributors, and popularity sampling may favour mainstream English-language films, so results may not generalise to other film industries.

\textbf{Reporting and use of AI.} We report the failed construction, because omitting it would overstate the method's success on new data. Generative AI was used in this project for \todo{describe the actual use}, and all claims and references were checked by the author.

\section{Conclusions and Reflections}
\label{sec:conclusion}

Replication was successful. The official code ran with one device fix, results fell within one standard deviation of the paper, and the Max--1-layer baseline matched closely. GIN also generalised to new data once the graphs carried real structure, beating the node-count baseline on every such dataset. Its theoretical advantage appeared clearly only on ENZYMES; on film ego-networks it was within noise, as on the original data. Higher expressive power therefore did not guarantee a measurable gain on every task.

The most useful lesson came from our own mistake. Mean density of exactly 1.000 was visible from the first dataset we built, and the accuracy drop looked like a finding about GIN until we tested the data itself. A simple baseline and a clique count would have caught it before any GPU time was spent, and both now run first. The same applies to the date filter, which failed silently, so a filter needs checking against its output. A further lesson is that a data source's terms need reading in full before collection. We also learned to measure before assuming: a short timing test showed the GPU was the slower choice for GraphSAGE on dense graphs.

\bibliography{references}

\end{document}
